% Hướng dẫn cho Văn phòng (VP). Dịch: pdflatex slides.tex (hai lần), sau scripts/test-all.sh
\documentclass[aspectratio=169,12pt]{beamer}
\usepackage{../pi-hd}
\newcommand\vaitro{Hướng dẫn cho Văn phòng}
\begin{document}

\bia{Văn phòng}{Nhận bài của tác giả, đưa vào hệ thống, giữ liên hệ của tác giả.}

\chu{Việc của Văn phòng}{%
\begin{buoc}
\item Nhận bài qua email, bưu điện hay trực tiếp.
\item \textbf{Thêm bài} trên trang: tác giả (mới hoặc đã có), mỗi bài một khối, ảnh kèm chỗ đặt.
\item Hệ thống cấp mã: tháng nhận, số thư mục kế tiếp, a, b, …
\item Liên hệ của tác giả chỉ Văn phòng, Phụ trách chuyên mục, Tổng biên tập, Quản trị xem — \textbf{không ghi vào đề hay lời giải}.
\item Khi cần hỏi tác giả (theo mục cần kiểm tra), Văn phòng liên lạc.
\end{buoc}}

\manhinh{them-bai}{Thêm bài: tác giả và tệp}
\manhinh[Ảnh lộ đáp án phải đặt ở lời giải — ảnh của đề được in kèm đề.]{them-bai-2}{Thêm bài: ảnh và lưu}
\manhinh{da-them}{Bài vừa thêm}

\chu{Khi có sự cố}{%
\begin{itemize}\setlength\itemsep{2mm}
\item \textbf{Bị nhắc „Chọn chỗ đặt cho ảnh"}: mỗi ảnh chọn hình của đề (in) hay minh hoạ lời giải (không in).
\item \textbf{Tệp .tex không tự tách}: dòng „Lời giải" phải đứng riêng; tách tay bằng cách chép vào hai ô.
\item \textbf{Tác giả đã có}: chọn trong danh sách Tác giả, đừng tạo trùng.
\end{itemize}}
\end{document}
